% Created 2026-01-22 jue 16:18
% Intended LaTeX compiler: pdflatex
\documentclass[a4paper]{article}
\usepackage[utf8]{inputenc}
\usepackage[T1]{fontenc}
\usepackage{graphicx}
\usepackage{longtable}
\usepackage{wrapfig}
\usepackage{rotating}
\usepackage[normalem]{ulem}
\usepackage{amsmath}
\usepackage{amssymb}
\usepackage{capt-of}
\usepackage{hyperref}
\usepackage{color}
\usepackage{listings}
\usepackage[spanish]{babel}
\usepackage[usenames,dvipsnames]{color} % Required for custom colors
\renewcommand{\ttdefault}{pcr} % MONOESPACIO CON NEGRIT
\usepackage{lastpage}
\usepackage{listings}
\usepackage{listingsutf8}
\renewcommand{\lstlistingname}{Listado}
\lstset{frame=single,inputencoding=utf8,basicstyle=\scriptsize\ttfamily,showstringspaces=false,numbers=none}
\definecolor{MyDarkGreen}{rgb}{0.0,0.4,0.0} % This is the color used for comments
\lstset{ breaklines=true, postbreak=\mbox{\textcolor{red}{$\hookrightarrow$}\space}, keywordstyle=\bfseries, keywordstyle=[1]\color{Blue}\bfseries,  keywordstyle=[2]\color{Purple}\bfseries,  keywordstyle=[3]\color{Blue}\underbar,   identifierstyle=,   commentstyle=\usefont{T1}{pcr}{m}{sl}\color{MyDarkGreen}\small,   stringstyle=\color{Purple},   showstringspaces=false,   tabsize=2,   morecomment=[l][\color{Blue}]{...} }
\lstset{literate=  {á}{{\'a}}1 {é}{{\'e}}1 {í}{{\'i}}1 {ó}{{\'o}}1 {ú}{{\'u}}1   {Á}{{\'A}}1 {É}{{\'E}}1 {Í}{{\'I}}1 {Ó}{{\'O}}1 {Ú}{{\'U}}1   {à}{{\`a}}1 {è}{{\`e}}1 {ì}{{\`i}}1 {ò}{{\`o}}1 {ù}{{\`u}}1   {À}{{\`A}}1 {È}{{\'E}}1 {Ì}{{\`I}}1 {Ò}{{\`O}}1 {Ù}{{\`U}}1   {ä}{{\"a}}1 {ë}{{\"e}}1 {ï}{{\"i}}1 {ö}{{\"o}}1 {ü}{{\"u}}1   {Ä}{{\"A}}1 {Ë}{{\"E}}1 {Ï}{{\"I}}1 {Ö}{{\"O}}1 {Ü}{{\"U}}1   {â}{{\^a}}1 {ê}{{\^e}}1 {î}{{\^i}}1 {ô}{{\^o}}1 {û}{{\^u}}1   {Â}{{\^A}}1 {Ê}{{\^E}}1 {Î}{{\^I}}1 {Ô}{{\^O}}1 {Û}{{\^U}}1   {œ}{{\oe}}1 {Œ}{{\OE}}1 {æ}{{\ae}}1 {Æ}{{\AE}}1 {ß}{{\ss}}1   {ű}{{\H{u}}}1 {Ű}{{\H{U}}}1 {ő}{{\H{o}}}1 {Ő}{{\H{O}}}1   {ç}{{\c c}}1 {Ç}{{\c C}}1 {ø}{{\o}}1 {å}{{\r a}}1 {Å}{{\r A}}1   {€}{{\euro}}1 {£}{{\pounds}}1 {«}{{\guillemotleft}}1   {»}{{\guillemotright}}1 {ñ}{{\~n}}1 {Ñ}{{\~N}}1 {¿}{{?`}}1 }
\usepackage{caption}
\usepackage{attachfile2}
\usepackage[margin=1.5cm,includeheadfoot,includehead,includefoot]{geometry}
\hypersetup{colorlinks,linkcolor=black}
\usepackage{fancyhdr}
\pagestyle{fancyplain}
\chead{}
\lhead{}
\rhead{}
\cfoot{}
\lfoot{\begin{footnotesize}alvaro.gonzalezsotillo@educa.madrid.org\end{footnotesize}}
\rfoot{\begin{footnotesize}\thepage / \pageref{LastPage}\end{footnotesize}}
\usepackage[skins]{tcolorbox}
\usepackage{multicol}
\usepackage{changepage} %ajdustwidth
\usepackage{fancybox}
\usepackage{attachfile2}
\lhead{Extraordinaria 2021 (es el \\lhead)}
\rhead{Administración y Gestión de Bases de Datos (es el \\rhead)}
\lhead{IAW}
\rhead{Práctica PHP (II)}
\usepackage{skak}
\usepackage{svg}
\usepackage{letltxmacro}
\LetLtxMacro{\originalincludegraphics}{\includegraphics}
\renewcommand{\includegraphics}[2][]{\IfFileExists{#2.pdf}{\originalincludegraphics[#1]{#2.pdf}}{\originalincludegraphics[#1]{#2}}}
\LetLtxMacro{\originalincludesvg}{\includesvg}
\renewcommand{\includesvg}[2][]{\IfFileExists{#2.pdf}{\originalincludegraphics[#1]{#2.pdf}}{\originalincludegraphics[#1]{#2.svg.pdf}}}
\usepackage{comment}
\excludecomment{NOTES}
\date{}
\title{Segunda práctica PHP}
\hypersetup{
 pdfauthor={Álvaro González Sotillo},
 pdftitle={Segunda práctica PHP},
 pdfkeywords={},
 pdfsubject={},
 pdfcreator={Emacs 30.2 (Org mode 9.7.11)}, 
 pdflang={Spanish}}
\begin{document}

\maketitle
\setcounter{tocdepth}{1}
\tableofcontents

\captionsetup{font=scriptsize}

\textattachfile{/datos-luks/datos-alvaro/apuntes-clase/implantacion-aplicaciones-web-asir2/apuntes/UT05/UT05-practica-ajedrez-php.org}{} \vspace{-1em}


\setlength{\parindent}{0em}
\setlength{\parskip}{1em}

\newtcolorbox{Aviso}[1][Aviso]{
  enhanced,
  colback=gray!5!white,
  colframe=gray!75!black,fonttitle=\bfseries,
  colbacktitle=gray!85!black,
  attach boxed title to top left={yshift=-2mm,xshift=2mm},
  title=#1
}

\newtcolorbox{cuadrito}[1][Ignorado]{
  %drop shadow southeast,
  enhanced jigsaw,
  colback=white,
}


\newcommand{\StudentData}{
  \begin{cuadrito}[1\textwidth]
    \vspace{0.3cm}
    \large{
      \textbf{Apellidos:} \hrulefill \\
      \textbf{Nombre:} \hrulefill \\
      \textbf{Fecha:} \hrulefill \hspace{1cm} \textbf{Usuario:} \hrulefill \\
    }
    \vspace{-0.2cm}
  \end{cuadrito}
}

\newcommand{\StudentDataDniFirma}{
  \begin{cuadrito}[1\textwidth]
    \vspace{0.3cm}
    \large{
      \textbf{Apellidos:} \hrulefill \\
      \textbf{Nombre:} \hrulefill  \\
      \textbf{Dni:} \hrulefill \hspace{8cm} \\
      \\
      \textbf{Firma:} \hrulefill \hspace{8cm} \\
    }
    \vspace{-0.2cm}
  \end{cuadrito}
}

\attachfile{../UT04//tiles/tiles.zip}
\section*{Objetivos de la práctica:}
\label{sec:org0000000}
En esta práctica se espera que el alumno:
\begin{itemize}
\item Se familiarice con el manejo de sesiones
\item Se familiarice con el uso de MySQL
\end{itemize}
\section*{Creación de tablas}
\label{sec:org0000003}
\begin{Aviso}[No se corregirá la práctica si...]
Estas son normas necesarias para poder corregir la práctica:
\begin{itemize}
\item La aplicación creará las tablas necesarias para su funcionamiento
\item Se puede comenzar la aplicación por cualquier página (no debe suponerse que siempre se lanzará \texttt{mitablero.php} antes de \texttt{administracion.php}, por ejemplo)
\item La información de conexión a la base de datos estará en un único fichero PHP. El resto de ficheros usarán esa información (\texttt{require()} o \texttt{include()})
\end{itemize}
\end{Aviso}
\section{Tablero personal (2 puntos)}
\label{sec:org0000006}
\begin{itemize}
\item En la página \texttt{mitablero.php}
\begin{itemize}
\item Se mostrará un tablero de ajedrez, inicialmente vacío, con indicaciones de filas y columnas (10\%)
\item Cada sesión de usuario tendrá un tablero distinto (40\%)
\item Habrá un formulario que permitirá introducir:
\begin{itemize}
\item Una posición del tablero (\texttt{a8}, \texttt{b3},\ldots{})
\item Seleccionar una ficha de ajedrez, o una casilla vacía, con un \texttt{<select>}
\item Al enviar el formulario, se anotará esa información en la base de datos y se mostrará el nuevo tablero (50\%)
\item Nota: no se tendrán en cuenta las normas del juego, pueden crearse cualquier disposición de fichas
\end{itemize}
\item Las fichas de ajedrez pueden ser caracteres unicode

\begin{html}
♔ ♕ ♖ ♗ ♘ ♙
♚ ♛ ♜ ♝ ♞ ♟
\end{html}
\begin{latex}
\BlackBishopOnWhite \BlackRookOnWhite \BlackKingOnWhite  \BlackKnightOnWhite \BlackPawnOnWhite \BlackQueenOnWhite
\WhiteBishopOnWhite \WhiteRookOnWhite \WhiteKingOnWhite  \WhiteKnightOnWhite \WhitePawnOnWhite \WhiteQueenOnWhite
\end{latex}
\begin{itemize}
\item Ver \url{https://en.wikipedia.org/wiki/Chess\_symbols\_in\_Unicode}
\end{itemize}
\end{itemize}
\end{itemize}

\begin{center}
\includesvg[width=0.4\textwidth]{media/chess}
\end{center}
\section{Tablero compartido (3 puntos)}
\label{sec:org0000009}
\begin{itemize}
\item Añade en \texttt{mitablero.php} un enlace para compartir el tablero
\begin{itemize}
\item El enlace apuntará a \texttt{tablerocompartido.php}, y añadirá un parámetro GET para saber qué tablero se comparte (10\%)
\end{itemize}
\item Una vez en \texttt{tablerocompartido.php}, la página se comportará igual que \texttt{mitablero.php}, pero modificando el tablero indicado en el parámetro GET, en vez del tablero asociado a la sesión
\begin{itemize}
\item Esto incluye el tablero, el formulario para cambiar el tablero y el enlace de compartir (90\%)
\end{itemize}
\item Si \texttt{mitablero.php} y \texttt{tablerocompartido.php} no comparten el código PHP y HTML principal (usando \texttt{require()} o \texttt{include()}), este ejercicio contará como máximo la mitad de la puntuación.
\end{itemize}
\section{Página de administración (4 puntos)}
\label{sec:org000000c}
\begin{itemize}
\item En la página \texttt{administracion.php}
\begin{itemize}
\item Se listarán todos los tableros, indicando su identificador
\item Cada tablero tendrá
\begin{itemize}
\item Un enlace a su página en \texttt{tablerocompartido.php} (10\%)
\item Un botón para borrarlo definitivamente (20\%)
\end{itemize}
\end{itemize}
\item Habrá un botón para crear un tablero compartido nuevo
\begin{itemize}
\item Con todas las fichas iniciales de un juego normal de ajedrez (50\%)
\item Al pulsarse, se irá al nuevo tablero compartido directamente (\href{https://stackoverflow.com/questions/768431/how-do-i-make-a-redirect-in-php}{ver \texttt{header( "Location: ...")}}) (20\%)
\end{itemize}
\end{itemize}
\section*{Normas de entrega}
\label{sec:org000000f}
\begin{itemize}
\item El código estará en el directorio \texttt{\textasciitilde{}/public-html/practica02/} del \emph{workspace} de Coder del alumno
\begin{itemize}
\item El profesor copiará las páginas a su propio LAMP para corregir la práctica
\item El profesor se encargará de cambiar los datos de acceso a MySQL.
\end{itemize}
\item Los nombres de fichero PHP y los nombres de funciones serán los especificados en los enunciados
\item Todos los formularios serán \textbf{GET}
\begin{itemize}
\item Si no se pueden ver los parámetros en la URL, el ejercicio no tendrá más del 50\% de la nota
\end{itemize}
\item Se recuerda que los alumnos deben guardar una copia de seguridad de todos sus trabajos, en el caso de que Coder no funcione
\item \textbf{No deben aparcer errores ni avisos de PHP}. Hay que tener en cuenta que se ejecutarán con opciones de \emph{debug}, como \texttt{display\_errors} a \texttt{true}
\begin{itemize}
\item Puede ser útil: \texttt{isset()}, \texttt{is\_numeric()}
\item Si hay errores o avisos, el ejercicio no tendrá más del 50\% de la nota
\end{itemize}
\item No se tendrán en cuenta:
\begin{itemize}
\item Errores debidos al reenvío de formularios por duplicado (por recarga de página)
\item Errores debidos a parámetros cambiados a mano
\end{itemize}
\end{itemize}
\section*{Resultados de aprendizaje}
\label{sec:org0000012}
Esta práctica contribuye a la calificación de los siguientes RA:
\begin{itemize}
\item RA 5: parte de "prácticas"
\item RA 6: parte de "prácticas"
\end{itemize}
\end{document}
